\section{Method and hierarchical extensions}
\label{sec:exec-method}

\subsection{Segment calculations and partition inference}

For a partition $\tau=(0=t_0<t_1<\cdots<t_k=n)$, define a local segment score
\[
  A_{ij}=M_{ij}\,c(i,j)\,h(j)^{\mathbbm 1\{j<n\}},
\]
where $c(i,j)$ is a segment-cohesion term and $h(j)$ is an optional prior factor for an interior boundary at coordinate $j$. Treating these as separate objects prevents a physical segment-length prior from being confused with a local boundary hazard.

The posterior weight of a partition with $k$ segments is proportional to
\[
  p(k)\prod_{r=1}^{k} A_{t_{r-1}+1,t_r}.
\]
Sum-product dynamic programming sums these weights over all admissible partitions. A max-sum version replaces summation by maximization and stores backpointers to recover the joint MAP partition. After all reachable segment scores have been tabulated, the worst-case recursion cost is $\mathcal O(k_{\max}n^2)$.

\begin{figure}[H]
  \centering
  \resizebox{0.96\textwidth}{!}{\begin{tikzpicture}[node distance=8mm and 10mm]
\node[primarybox,text width=31mm] (data) {Ordered observations\\$u,y,d,a$};
\node[secondarybox,right=of data,text width=35mm] (segment) {Observation-model calculation\\$\log A^{(r)}_{ij}$};
\node[primarybox,above right=7mm and 13mm of segment,text width=32mm] (sum) {Sum-product recursion\\marginal likelihoods and posterior marginals};
\node[primarybox,below right=7mm and 13mm of segment,text width=32mm] (max) {Max-sum recursion\\joint MAP partition};
\node[successbox,right=13mm of sum,text width=31mm] (post) {$p(k\mid y)$, changepoint probabilities, Bayes curves};
\node[successbox,right=13mm of max,text width=31mm] (map) {Joint MAP changepoints and fitted segments};
\draw[arrPrimary] (data)--(segment);
\draw[arrPrimary] (segment)--(sum);
\draw[arrPrimary] (segment)--(max);
\draw[arrPrimary] (sum)--(post);
\draw[arrPrimary] (max)--(map);
\node[draw=Rule,dashed,rounded corners,fit=(sum)(max),inner sep=4mm,label=above:{\small dynamic programming over ordered partitions}] {};
\end{tikzpicture}
}
  \caption{Statistical structure of BayesBreak. Observation-family-specific segment calculations feed two distinct dynamic-programming recursions: sum-product for posterior marginalization and max-sum with backtracking for the joint MAP partition.}
  \label{fig:exec-segment-partition}
\end{figure}

\subsection{Irregular designs}

Irregular coordinates $x_1<\cdots<x_n$ enter through the prior on partitions rather than through an index-only approximation. Segment cohesion may depend on physical span, while a boundary hazard may depend on a local design gap. The two factors have different probability interpretations and must be specified separately. The existing archived length-aware well-log fit is a real computation, but its former Poisson-gap interpretation is not retained until the implemented prior matches that derivation.

\subsection{Multiple sequences and known groups}

For aligned sequences $s=1,\ldots,S$ that are conditionally independent given a common partition, the segment marginal likelihood is
\[
  M^{\mathrm{shared}}_{ij}=\prod_{s=1}^{S}M^{(s)}_{ij}.
\]
The same partition recursion can then be applied to the pooled segment scores. Known groups use the same construction within each group. The shared-boundary model is appropriate only when the candidate coordinates are aligned and a common partition is substantively plausible.

\subsection{Finite latent groups}

When group membership is unknown, the current method optimizes a finite criterion over group weights, sequence responsibilities, and segmentation templates. A Jensen minorization produces normalized auxiliary weights, and each template update is solved by a responsibility-weighted max-sum recursion. This is a minorization--maximization procedure for the stated criterion. Its component scores are not normalized sampling densities, so the method is not described as a conventional finite-mixture likelihood and no mixture-identifiability result is claimed.

\begin{figure}[H]
  \centering
  \resizebox{0.72\textwidth}{!}{\begin{tikzpicture}[node distance=7mm and 8mm]
\node[neutralbox,text width=28mm] (part) {Common ordered partition\\$\tau=(t_0,\ldots,t_k)$};
\node[secondarybox,below left=9mm and 5mm of part,text width=25mm] (s1) {Sequence 1\\segment scores $M^{(1)}_{ij}$};
\node[secondarybox,below=9mm of part,text width=25mm] (s2) {Sequence $s$\\segment scores $M^{(s)}_{ij}$};
\node[secondarybox,below right=9mm and 5mm of part,text width=25mm] (sS) {Sequence $S$\\segment scores $M^{(S)}_{ij}$};
\node[primarybox,below=13mm of s2,text width=40mm] (pool) {Pooled segment score\\$\log M^{\rm shared}_{ij}=\sum_{s=1}^{S}\log M^{(s)}_{ij}$};
\draw[arrPrimary] (part)--(s1);
\draw[arrPrimary] (part)--(s2);
\draw[arrPrimary] (part)--(sS);
\draw[arrSecondary] (s1)--(pool);
\draw[arrSecondary] (s2)--(pool);
\draw[arrSecondary] (sS)--(pool);
\end{tikzpicture}
}
  \caption*{(a) Common partition across conditionally independent sequences.}
  \vspace{7pt}
  \resizebox{0.94\textwidth}{!}{\begin{tikzpicture}[node distance=8mm and 9mm]
\node[secondarybox,text width=31mm] (scores) {Sequence-specific template scores\\$S_s(\tau_g)>0$};
\node[primarybox,right=of scores,text width=29mm] (resp) {Auxiliary allocation weights\\$r_{sg}$};
\node[primarybox,right=of resp,text width=32mm] (update) {Alternating updates\\$\pi_g$ and $\tau_g$};
\node[successbox,right=of update,text width=31mm] (obj) {$\mathcal F=\sum_s\log\sum_g\pi_gS_s(\tau_g)$};
\draw[arrPrimary] (scores)--(resp);
\draw[arrPrimary] (resp)--(update);
\draw[arrPrimary] (update)--(obj);
\draw[arrSecondary] (obj.north) |- ++(0,7mm) -| node[pos=0.52,above,font=\scriptsize]{next iteration} (scores.north);
\node[neutralbox,below=10mm of resp,text width=101mm] {The updates maximize a Jensen minorizer of the stated finite latent-group criterion. The criterion is not interpreted as a normalized finite-mixture likelihood.};
\end{tikzpicture}
}
  \caption*{(b) Finite latent-group allocation and template updates.}
  \caption{Hierarchical extensions. The shared-boundary model pools sequence-specific segment marginal likelihoods. The latent-group model alternates responsibility and template updates for the explicitly stated finite criterion.}
  \label{fig:exec-hierarchy}
\end{figure}
